\subsection{Adaptive Exploration Mechanism}
\label{sec:adaptive_epsilon_decay}

In the proposed event-driven scheduling framework, agent interactions with the
environment occur exclusively at \emph{decision points} triggered by execution
events such as job arrivals and operation completions
(Section~\ref{sec:discrete_event}). Consequently, the number of decision points
within an episode is neither fixed nor known a priori and varies stochastically
across episodes.

This property fundamentally affects how exploration should be scheduled.
Exploration must be controlled with respect to \emph{decision experience}, not
elapsed simulation time or episode count.

\subsubsection{Why Standard Decay Mechanisms Fail}

\paragraph{Standard time-based $\epsilon$ decay.}
In conventional reinforcement learning, the exploration rate $\epsilon$ is often
decayed as a function of elapsed time or step index. A common linear decay
formulation is
\begin{equation}
\epsilon(t)
=
\max\!\left(
\epsilon_{\min},
\;
\epsilon_0 - \frac{t}{T_{\mathrm{anneal}}}
(\epsilon_0 - \epsilon_{\min})
\right),
\label{eq:time_based_epsilon}
\end{equation}
where $t$ denotes elapsed simulation time (or episode step), $T_{\mathrm{anneal}}$
is a fixed annealing horizon, and $\epsilon_0,\epsilon_{\min}$ are the initial and
minimum exploration rates, respectively.

This formulation implicitly assumes that the number of meaningful agent–environment
interactions grows \emph{proportionally with} $t$. In discrete-time environments
with fixed step intervals, this assumption holds: each time step corresponds to
exactly one decision opportunity. However, in the considered event-driven setting,
this assumption is violated.

\paragraph{Episode-based $\epsilon$ decay.}
Alternatively, one might decay exploration as a function of episode count:
\begin{equation}
\epsilon(e)
=
\max\!\left(
\epsilon_{\min},
\;
\epsilon_0 - \frac{e}{E_{\mathrm{anneal}}}
(\epsilon_0 - \epsilon_{\min})
\right),
\label{eq:episode_based_epsilon}
\end{equation}
where $e$ is the current episode index and $E_{\mathrm{anneal}}$ is the number of
episodes over which exploration should be reduced.

This approach is commonly used in multi-agent reinforcement learning and assumes
that each episode provides approximately equal learning experience. For example,
setting $E_{\mathrm{anneal}} = 600$ with $1000$ total episodes implies that
exploration should dominate the first $60\%$ of training episodes, allowing
agents to accumulate sufficient experience before transitioning to exploitation.

However, this assumption fails when the number of decision points per episode
varies stochastically.

\paragraph{Formal problem statement.}
Let $D_e$ denote the number of decision points in episode $e$. In event-driven
scheduling with stochastic job arrivals and heterogeneous job structures,
$D_e$ is a random variable:
\begin{equation}
D_e \sim \mathcal{D}(\lambda, \mu, \sigma),
\label{eq:decision_points_distribution}
\end{equation}
where $\lambda$ represents arrival rate, $\mu$ represents mean operations per job,
and $\sigma$ captures other sources of stochasticity (processing times, machine
breakdowns, etc.).

Crucially, $D_e$ can vary significantly across episodes. In the considered
experimental setting (Section~\ref{sec:experiments}), observed decision point
counts range from $D_{\min}=37$ to $D_{\max}=51$ per episode, with
$\bar{D}=41.66$ and $\sigma_D=2.70$.

\paragraph{Experience accumulation imbalance.}
Under episode-based decay (Eq.~\ref{eq:episode_based_epsilon}), the \emph{total
number of exploratory decisions} accumulated during the annealing phase is:
\begin{equation}
N_{\mathrm{explore}}^{\mathrm{total}}
=
\sum_{e=1}^{E_{\mathrm{anneal}}} \epsilon(e) \cdot D_e,
\label{eq:total_explore_episode}
\end{equation}
where $\epsilon(e) \cdot D_e$ approximates the expected number of exploratory
actions in episode $e$ (assuming $\epsilon$-greedy selection).

Because $D_e$ varies stochastically, two training runs with identical
hyperparameters can accumulate vastly different amounts of exploration experience:
\begin{align}
N_{\mathrm{explore}}^{\mathrm{Run~A}}
&= \sum_{e=1}^{600} \epsilon(e) \cdot D_e^{(A)},
\label{eq:explore_run_a}
\\
N_{\mathrm{explore}}^{\mathrm{Run~B}}
&= \sum_{e=1}^{600} \epsilon(e) \cdot D_e^{(B)},
\label{eq:explore_run_b}
\end{align}
where $D_e^{(A)}$ and $D_e^{(B)}$ are realizations from the same distribution.

\paragraph{Numerical example.}
Consider $E_{\mathrm{anneal}}=600$ episodes with linear decay from
$\epsilon_0=1.0$ to $\epsilon_{\min}=0.05$. Assume two scenarios:

\begin{itemize}
\item \textbf{Scenario A (low event density):}  
First $600$ episodes happen to have low decision point counts, with
$D_e^{(A)} \approx 38$ on average.

\item \textbf{Scenario B (high event density):}  
First $600$ episodes happen to have high decision point counts, with
$D_e^{(B)} \approx 45$ on average.
\end{itemize}

The average exploration rate over the annealing phase is approximately:
\begin{equation}
\bar{\epsilon}_{\mathrm{anneal}}
=
\frac{1}{E_{\mathrm{anneal}}} \sum_{e=1}^{E_{\mathrm{anneal}}} \epsilon(e)
\approx
\frac{\epsilon_0 + \epsilon_{\min}}{2}
=
0.525.
\label{eq:avg_epsilon}
\end{equation}

Therefore, the expected number of exploratory decisions is:
\begin{align}
\mathbb{E}[N_{\mathrm{explore}}^{A}]
&\approx 0.525 \times 600 \times 38
= 11{,}970,
\label{eq:expected_explore_a}
\\
\mathbb{E}[N_{\mathrm{explore}}^{B}]
&\approx 0.525 \times 600 \times 45
= 14{,}175.
\label{eq:expected_explore_b}
\end{align}

This represents a \textbf{18.4\% difference in exploration experience} despite
identical episode counts and decay parameters. In Scenario~A, agents prematurely
transition to exploitation after accumulating only $11{,}970$ exploratory
decisions, whereas Scenario~B agents receive $14{,}175$ decisions.

More critically, the \emph{remaining} $400$ episodes (intended for exploitation)
may contain more decision points than the first $600$ episodes. If episodes
$601$–$1000$ happen to have higher event density ($\bar{D}=48$), they contain:
\begin{equation}
D_{\mathrm{remaining}}
=
400 \times 48
=
19{,}200
\text{ decision points}.
\label{eq:remaining_decisions}
\end{equation}

This exceeds the exploration phase in Scenario~A, meaning that more decision
opportunities are wasted during exploitation than were used during exploration.

\paragraph{Time-based decay fails similarly.}
Time-based decay (Eq.~\ref{eq:time_based_epsilon}) suffers from the same problem.
Consider two episodes with identical duration $T=40$ but different decision point
densities:

\begin{itemize}
\item \textbf{Episode A:} $D_A=10$ decision points occur at
$t \in \{2.5, 5.0, 8.3, 12.1, 18.5, 24.0, 28.7, 31.2, 35.6, 38.9\}$.

\item \textbf{Episode B:} $D_B=40$ decision points occur uniformly at
$t \in \{1, 2, 3, \ldots, 40\}$.
\end{itemize}

Setting $T_{\mathrm{anneal}}=T=40$ means exploration is fully suppressed at
$t=40$. With linear decay, $\epsilon(t) \approx 0.4$ at $t=24$ (the $60\%$ point).

At this time:
\begin{itemize}
\item Episode~A has encountered only $d_A(24)=6$ decision points.
\item Episode~B has encountered $d_B(24)=24$ decision points.
\end{itemize}

Thus, Episode~A agents have explored only $6$ decisions before $\epsilon$ drops
below $0.4$, while Episode~B agents have explored $24$ decisions. The mapping
\begin{equation}
d(t) \neq \alpha t
\label{eq:decision_time_mismatch}
\end{equation}
is \emph{non-linear} and \emph{stochastic}, violating the proportionality
assumption underlying time-based decay.

\subsubsection{Decision-Point-Based Exploration Decay}

\paragraph{Proposed formulation.}
The proposed framework defines exploration as a direct function of the
\emph{cumulative decision point index} $d$:
\begin{equation}
\epsilon_d
=
\max\!\left(
\epsilon_{\min},
\;
\epsilon_0 - \frac{d}{D_{\mathrm{anneal}}}
(\epsilon_0 - \epsilon_{\min})
\right),
\label{eq:decision_based_epsilon}
\end{equation}
where $d$ counts the total number of decision points encountered across all
episodes so far, and $D_{\mathrm{anneal}}$ specifies the number of decision points
over which exploration is reduced.

Critically, $d$ is \emph{incremented at each decision point}, not at each time
step or episode boundary. This ensures that exploration decay is directly coupled
to the interaction structure of the environment.

\paragraph{Guaranteed exploration coverage.}
Under decision-based decay, the total number of exploratory decisions during the
annealing phase is deterministic (up to $\epsilon$-greedy stochasticity):
\begin{equation}
N_{\mathrm{explore}}^{\mathrm{total}}
\approx
\int_0^{D_{\mathrm{anneal}}} \epsilon_d \, \mathrm{d}d
=
\frac{D_{\mathrm{anneal}}(\epsilon_0 + \epsilon_{\min})}{2}.
\label{eq:guaranteed_explore}
\end{equation}

For example, setting $D_{\mathrm{anneal}}=25{,}000$ with $\epsilon_0=1.0$ and
$\epsilon_{\min}=0.05$ guarantees approximately:
\begin{equation}
N_{\mathrm{explore}}
\approx
\frac{25{,}000 \times (1.0 + 0.05)}{2}
=
13{,}125
\text{ exploratory decisions},
\label{eq:guaranteed_explore_example}
\end{equation}
\emph{regardless of how these decision points are distributed across episodes}.

This property ensures that:
\begin{enumerate}
\item Agents accumulate \emph{consistent exploration experience} across different
training runs, even when episode-level stochasticity varies.

\item Exploration does not terminate prematurely in episodes with low event
density, nor does it persist excessively in high-density episodes.

\item The fraction of decision points allocated to exploration versus exploitation
is explicitly controlled via $D_{\mathrm{anneal}}$.
\end{enumerate}

\paragraph{Episode-level comparison.}
To compare episode-based and decision-based decay, consider the same $1000$-episode
training scenario:

\begin{itemize}
\item \textbf{Episode-based decay:} $E_{\mathrm{anneal}}=600$ episodes.
\item \textbf{Decision-based decay:} $D_{\mathrm{anneal}}=25{,}000$ decision points.
\end{itemize}

Assume decision points per episode vary as $D_e \sim \mathcal{N}(41.66, 2.70^2)$
(observed distribution in experiments). The total decision point count across
$1000$ episodes is:
\begin{equation}
D_{\mathrm{total}}
=
1000 \times 41.66
=
41{,}660.
\label{eq:total_decisions}
\end{equation}

Under episode-based decay, the exploration phase covers approximately:
\begin{equation}
D_{\mathrm{explore}}^{\mathrm{episode}}
=
600 \times 41.66
=
24{,}996
\text{ decision points (expected)}.
\label{eq:episode_explore_decisions}
\end{equation}

However, due to stochasticity, the actual count can range from:
\begin{equation}
D_{\mathrm{explore}}^{\mathrm{episode}}
\in
[600 \times 37, \; 600 \times 51]
=
[22{,}200, \; 30{,}600].
\label{eq:episode_explore_range}
\end{equation}

This $38\%$ variation in exploration experience is eliminated under decision-based
decay, where exactly $D_{\mathrm{anneal}}=25{,}000$ decision points are explored.

\paragraph{Implications for learning.}
In dynamic job shop scheduling, informative learning signals arise from resource
conflicts, congestion patterns, and asynchronous interactions
(Section~\ref{sec:parallel_conflict}). These signals are encountered at decision
points rather than at uniform time intervals or episode boundaries.

Decision-point-based exploration decay therefore:
\begin{itemize}
\item \textbf{Aligns exploration with interaction structure:} Decay proceeds at
the rate agents actually interact with the environment.

\item \textbf{Prevents premature convergence:} Low-density episodes do not cause
premature exploitation.

\item \textbf{Ensures robust learning:} Consistent exploration experience across
runs with different stochastic realizations.

\item \textbf{Enables principled hyperparameter tuning:} $D_{\mathrm{anneal}}$ has
a clear semantic meaning (number of exploratory decisions), unlike
$E_{\mathrm{anneal}}$ (number of episodes, whose decision content varies).
\end{itemize}

This mechanism is essential for achieving stable and reproducible learning in
stochastic, event-driven multi-agent scheduling environments.
